% Requer \usepackage{xcolor} no preâmbulo do documento principal.
\begingroup
\color{blue}
\section{Discussão}
\label{sec:discussao}

Na base completa, o Extra Trees (ET) com seis sensores MQ e três variáveis
ambientais identificou corretamente 9 de 10 vasos, 31 de 36 coletas e
55.136 de 63.492 leituras (86,84\%) na validação que reservou um vaso por
rodada. A separação entre os grupos mostra que, nas condições avaliadas, os
sinais registrados continham informação sobre a atribuição dos vasos ao
grupo inoculado ou ao controle. Ela não identifica, por si, a fonte dessa
informação: os rótulos representam a inoculação planejada, sem confirmação
nematológica individual, e a matriz de sensores responde a misturas gasosas
sem identificar compostos. Tampouco demonstra reconhecimento específico de
\textit{H. glycines} ou desempenho no diagnóstico de infestação natural.
As leituras de um mesmo vaso são repetidas e correlacionadas; a unidade
biológica independente é o vaso, não cada uma das 63.492 leituras. Com
apenas cinco vasos por grupo, o intervalo binomial exato de 95\% para o
acerto de 9/10 vasos é amplo (55,50--99,75\%).

O contraste entre conjuntos de entradas restringe a interpretação do
sinal dos MQ. O ET com somente índice de umidade do solo, temperatura e
pressão acertou 84,71\% das leituras, 29 coletas e 10 vasos; com somente os
seis MQ, acertou 78,81\% das leituras, 28 coletas e 9 vasos. A inclusão dos
MQ elevou a acurácia por leitura em 2,13 pontos percentuais em relação ao
modelo ambiental, mas reduziu de 10 para 9 o número de vasos corretos.
A diferença média de acurácia por leitura dentro de vaso foi de 2,80
pontos percentuais, e o contraste exploratório por inversões de sinal não
forneceu evidência de ganho estável ($p=0{,}348$). Esse resultado não prova
que os MQ sejam irrelevantes; mostra que sua contribuição adicional ainda
não foi estabelecida neste conjunto de vasos. O índice de solo acompanha
a irrigação e o período de coleta; temperatura e pressão também podem
carregar informação compartilhada pelas medições de uma sessão. Além
disso, essas condições podem alterar o \textit{headspace} e a resposta dos
MQ. Assim, a classificação baseada em ambiente pode explorar diferenças
associadas ao calendário e ao procedimento de aquisição, inclusive quando
elas não são específicas do nematoide. Na análise de importância por
permutação do teste por coleta, a pressão teve grande contribuição
preditiva, sobretudo na TCN, mas essa medida não distingue um efeito
biológico da dinâmica da câmara e da operação do vácuo.

A dificuldade de generalização temporal reforça essa cautela. No ET com
nove entradas, a validação por vaso acertou 31/36 coletas, ao passo que a
avaliação que reservou uma data inteira acertou 10/36 com a faixa estrita
de pressão e 13/36 sem ela. As duas avaliações respondem a perguntas
distintas: a primeira testa novos vasos em datas já representadas no
treinamento; a segunda testa datas novas, mas pode manter os mesmos vasos
nos dois conjuntos. A queda é compatível com dependência de condições
temporais, ambientais ou operacionais e não permite determinar qual delas
responde pelo efeito. Na transferência entre os períodos de 11--13 e
18--19 de maio, o ET com nove entradas acertou 9/19 coletas no sentido
anterior--posterior e 6/17 no inverso. A irrigação entre esses períodos
coincidiu com a mudança do índice de solo, mas também mudaram as datas e
possivelmente outras condições de aquisição; essas comparações não isolam
o efeito causal da irrigação.

A posição relativa dos modelos variou entre os recortes exploratórios do
índice operacional do solo. No recorte denominado solo seco, a
TCN apresentou os melhores resultados, alcançando 98,47\% de acurácia,
98,03\% de acurácia balanceada e 98,38\% de F1 macro. No solo úmido, o
Extra Trees foi superior, com 74,46\%, 62,96\% e 63,33\%, respectivamente.

Ao considerar conjuntamente todas as leituras coletadas a vácuo, o Extra
Trees atingiu 90,83\% de acurácia e a TCN atingiu 89,91\%. Essa diferença
global é pequena, mas oculta comportamentos distintos: a TCN foi mais
eficaz no recorte seco e falhou na identificação da classe controle no recorte
úmido, enquanto o Extra Trees apresentou menor variação das métricas entre os
dois recortes.

Assim, os resultados não indicam superioridade geral de um único modelo.
A TCN apresentou métricas mais altas no recorte seco; o Extra Trees manteve
métricas mais altas no recorte úmido. Como os recortes correspondem a datas
distintas e houve irrigação entre esses períodos, essa comparação não identifica o
efeito da água sobre nenhum dos modelos.

Essa comparação pertence à partição histórica por coleta, na qual os
mesmos vasos puderam aparecer em treinamento e teste em dias distintos.
No recorte úmido, a TCN não reconheceu nenhuma leitura controle; sua
acurácia simples, influenciada pelo desbalanceamento, superestima a
capacidade de distinguir os dois grupos. O teste úmido reuniu somente
quatro coletas, e a distinção entre solo seco e úmido se baseou em um
índice operacional, sem medida calibrada da água em diferentes
profundidades. Há explicações concorrentes para a queda: a água pode
modificar a emissão, a retenção e o transporte dos voláteis no solo;
pode alterar a fisiologia da soja e da microbiota; o vapor de água pode
modificar a resposta dos sensores; e o protocolo, a composição das
classes ou a deriva instrumental podem diferir entre datas
\cite{asensio2007,cellini2017}. Como esses fatores variaram conjuntamente,
os dados não permitem escolher uma causa nem afirmar que a água reduziu
um sinal específico de \textit{H. glycines}.

A plausibilidade biológica da detecção por voláteis encontra apoio em
\cite{constantino2021}, que identificaram por GC--MS diferenças nas emissões
foliares de soja inoculada com \textit{H. glycines} antes de sintomas
evidentes. Nesse experimento, solo e raízes foram isolados durante a
coleta; portanto, seus compostos candidatos não demonstram que a mesma
assinatura esteja presente, ou seja detectada pelos MQ, no
\textit{headspace} integrado do sistema solo--planta estudado aqui. A
comparação com aplicações agrícolas de narizes eletrônicos também exige
atenção à especificidade. \cite{herrmann2024} distinguiram estresse hídrico
em soja com um nariz eletrônico e observaram variação dos sinais entre
períodos do dia. Logo, um padrão gasoso discriminatório pode refletir
estado hídrico ou horário, além da condição fitossanitária. O presente
resultado apoia a investigação de uma ferramenta de triagem em condições
controladas; sua utilidade agrícola dependerá de manter a discriminação
frente a esses fatores e a outros estresses, com confirmação independente
dos casos.

O alcance das estimativas também é limitado pelo desenho e pelas decisões
analíticas. Não há, nesta base, confirmação nematológica individual nem
registro de que os controles receberam o mesmo veículo e procedimento de
inoculação sem juvenis. Posição dos vasos e ordem de medição não foram
documentadas para esta análise, impedindo descartar efeitos de posição,
sequência ou deriva. A busca de subconjuntos de MQ utilizou os mesmos
vasos empregados na avaliação; por isso, o melhor resultado dessa busca
pode ser otimista. Cinco coletas foram excluídas após inspeção de erros e
sinais, sem comprovação de falha instrumental em todas elas. A acurácia
de 97,92\% obtida após essas exclusões é exploratória, não uma estimativa
independente do desempenho. Nenhum modelo foi validado em vasos, campanhas
ou locais externos ao experimento, e a validação por data já indica
instabilidade sob mudança de condições.

Os próximos ensaios devem separar as explicações concorrentes em um
desenho fatorial com mais vasos independentes: inoculação confirmada e
controle com inoculação simulada, cruzados com níveis de água do solo
medidos por método calibrado, em todas as datas. A posição dos vasos e a
ordem de amostragem devem ser aleatorizadas, enquanto horário, temperatura,
umidade do gás, pressão, vazão e fases da coleta são registrados. Controles
com outros nematoides, doenças, microbiotas e estresses abióticos testarão
se a resposta é específica de \textit{H. glycines}; diferentes cargas do
nematoide e tempos após a inoculação testarão o limite de detecção. A
coleta pareada de GC--MS e dos sinais MQ, inclusive em compartimentos
foliar e solo--raiz separados, poderá verificar quais compostos chegam à
matriz e se explicam a classificação. Por fim, seleção de entradas,
ajuste do modelo e definição de limiares devem ocorrer apenas nos dados
de treinamento, seguidos de avaliação cega em novos vasos e campanhas e,
depois, em amostras de campo com diagnóstico nematológico de referência.
Sensibilidade, especificidade e erros por vaso nessas condições determinarão
se o sistema é útil para triagem prática.
\endgroup
